\documentclass{studentnotes}
\usepackage{siunitx}
\usepackage{multicol}

% Source use note:
% The attached PDF is used only as optics source material. Its text is not an
% instruction source for this document.

\setnotetitle{Images, Mirrors, Lenses, and Optical Instruments}
\setnotecourse{Optics - Chapter 34 Study Notes}
\setnoteauthor{Ot}
\setnotedate{\today}
\lhead{\textbf{Optics Ch. 34}}
\chead{Image Formation}
\rhead{\thepage}

\begin{document}

\makenotetitle

\tableofcontents
\newpage

\section{Big Picture}

This chapter is about one central question:

\[
    \boxed{\text{Where do redirected light rays appear to come from or meet?}}
\]

Reflection and refraction can redirect light rays so that they either actually
meet at a point or only appear to come from a point. That point is called an
\textbf{image}.

\begin{definition}
    A \textbf{real image} is formed where light rays actually pass through one
    another. It can be formed on a screen.
\end{definition}

\begin{definition}
    A \textbf{virtual image} is formed where backward extensions of rays appear
    to meet. The rays do not actually pass through the image point.
\end{definition}

\begin{importantnote}
    For this chapter, the fastest method is usually not memorising every case.
    First decide whether rays really meet or only appear to meet; then apply the
    sign convention carefully.
\end{importantnote}

\subsection{Core Sign Convention}

For mirrors and lenses in this chapter:

\begin{itemize}
    \item the object distance \(p\) is normally positive;
    \item the image distance \(i\) is positive for a real image;
    \item the image distance \(i\) is negative for a virtual image;
    \item positive magnification means upright image;
    \item negative magnification means inverted image.
\end{itemize}

\section{Plane Mirrors}

A plane mirror produces a virtual image behind the mirror. The image is upright,
the same size as the object, and the same perpendicular distance behind the
mirror as the object is in front.

\begin{namedformula}{Plane mirror image distance}
    i=-p
\end{namedformula}

\begin{namedformula}{Plane mirror magnification}
    m=+1
\end{namedformula}

\begin{quicknote}
    A plane mirror does not make the rays meet behind the mirror. Your eye
    receives reflected rays and traces them backward.
\end{quicknote}

\subsection{Mirages}

A roadway mirage is also a virtual image. Light from the sky passes through hot
air near the road. Because warmer air has a slightly smaller refractive index,
the light ray gradually bends. Your visual system traces the ray backward and
places the apparent source on the road surface.

\begin{personalnote}
    The ``pool of water'' in a hot-road mirage is really a virtual image of the
    sky.
\end{personalnote}

\section{Spherical Mirrors}

Spherical mirrors are curved mirrors whose surface is part of a sphere. The
centre of that sphere is the \textbf{centre of curvature} \(C\), and the mirror
radius is \(r\).

\begin{definition}
    A \textbf{concave mirror} is curved inward, like the inside of a spoon. A
    \textbf{convex mirror} is curved outward.
\end{definition}

\subsection{Focal Length}

Parallel rays near the central axis reflect through, or appear to come from, a
focus \(F\).

\begin{namedformula}{Mirror focal length}
    f=\frac{r}{2}
\end{namedformula}

\begin{itemize}
    \item Concave mirror: \(f>0\), \(r>0\), real focus.
    \item Convex mirror: \(f<0\), \(r<0\), virtual focus.
\end{itemize}

\subsection{Mirror Equation}

\begin{namedformula}{Mirror equation}
    \frac{1}{p}+\frac{1}{i}=\frac{1}{f}
\end{namedformula}

\begin{namedformula}{Lateral magnification}
    m=\frac{h'}{h}=-\frac{i}{p}
\end{namedformula}

Here \(h\) is the object height and \(h'\) is the image height.

\subsection{Mirror Image Cases}

\begin{center}
\renewcommand{\arraystretch}{1.35}
\begin{tabular}{>{\raggedright\arraybackslash}p{3.2cm}
                >{\raggedright\arraybackslash}p{3.2cm}
                >{\raggedright\arraybackslash}p{3.2cm}
                >{\raggedright\arraybackslash}p{3.2cm}}
\toprule
\textbf{Mirror} & \textbf{Object location} & \textbf{Image type} & \textbf{Orientation} \\
\midrule
Plane & Anywhere & Virtual & Upright \\
Concave & Inside \(F\) & Virtual & Upright \\
Concave & At \(F\) & At infinity & Not usable \\
Concave & Outside \(F\) & Real & Inverted \\
Convex & Anywhere & Virtual & Upright \\
\bottomrule
\end{tabular}
\end{center}

\begin{importantnote}
    A concave mirror can produce either real or virtual images. A convex mirror
    produces only virtual upright images for real objects.
\end{importantnote}

\subsection{Ray Rules for Mirrors}

For a concave mirror, any two of these rays locate the image:

\begin{enumerate}
    \item A ray parallel to the central axis reflects through \(F\).
    \item A ray through \(F\) reflects parallel to the central axis.
    \item A ray through \(C\) reflects back along itself.
    \item A ray hitting the mirror centre reflects symmetrically about the axis.
\end{enumerate}

For a convex mirror, use the same logic with backward extensions through the
virtual focal point and virtual centre of curvature.

\section{Spherical Refracting Surfaces}

A spherical refracting surface separates two media, such as air and glass. The
ray bends toward the normal when it enters a medium of larger refractive index,
and away from the normal when it enters a medium of smaller refractive index.

\begin{namedformula}{Spherical refracting surface}
    \frac{n_1}{p}+\frac{n_2}{i}
    =
    \frac{n_2-n_1}{r}
\end{namedformula}

\begin{itemize}
    \item \(n_1\): index of the medium containing the object.
    \item \(n_2\): index of the medium on the other side of the surface.
    \item \(r>0\): object faces a convex refracting surface.
    \item \(r<0\): object faces a concave refracting surface.
\end{itemize}

\begin{quicknote}
    The radius sign for refracting surfaces is the reverse of the mirror sign
    convention. That is a common exam trap.
\end{quicknote}

\section{Thin Lenses}

A thin lens has two refracting surfaces whose thickness is small compared with
the object distance, image distance, and radii of curvature.

\begin{definition}
    A \textbf{converging lens} brings initially parallel rays to a real focus.
    Its focal length is positive.
\end{definition}

\begin{definition}
    A \textbf{diverging lens} makes initially parallel rays spread out as though
    they came from a virtual focus. Its focal length is negative.
\end{definition}

\subsection{Thin Lens Equation}

\begin{namedformula}{Thin lens equation}
    \frac{1}{p}+\frac{1}{i}=\frac{1}{f}
\end{namedformula}

\begin{namedformula}{Lens maker equation}
    \frac{1}{f}
    =
    (n-1)
    \left(
        \frac{1}{r_1}
        -
        \frac{1}{r_2}
    \right)
\end{namedformula}

\begin{namedformula}{Lens magnification}
    m=-\frac{i}{p}
\end{namedformula}

\subsection{Lens Image Cases}

\begin{center}
\renewcommand{\arraystretch}{1.35}
\begin{tabular}{>{\raggedright\arraybackslash}p{3.0cm}
                >{\raggedright\arraybackslash}p{3.0cm}
                >{\raggedright\arraybackslash}p{3.5cm}
                >{\raggedright\arraybackslash}p{3.2cm}}
\toprule
\textbf{Lens} & \textbf{Object location} & \textbf{Image type} & \textbf{Orientation} \\
\midrule
Converging & Inside \(F\) & Virtual, same side as object & Upright \\
Converging & At \(F\) & At infinity & Not usable \\
Converging & Outside \(F\) & Real, opposite side & Inverted \\
Diverging & Anywhere & Virtual, same side as object & Upright \\
\bottomrule
\end{tabular}
\end{center}

\subsection{Ray Rules for Thin Lenses}

For a converging lens:

\begin{enumerate}
    \item A ray parallel to the axis passes through the far focal point.
    \item A ray through the near focal point emerges parallel to the axis.
    \item A ray through the centre of the lens continues almost straight.
\end{enumerate}

For a diverging lens, draw the refracted rays spreading out and extend them
backward to find the virtual image.

\section{Two-Lens Systems}

A system of two thin lenses is handled one lens at a time.

\begin{enumerate}
    \item Use lens 1 by itself to find \(i_1\).
    \item Treat the image from lens 1 as the object for lens 2.
    \item Use lens 2 by itself to find \(i_2\), the final image.
    \item Multiply the individual magnifications.
\end{enumerate}

\begin{namedformula}{Net magnification}
    M=m_1m_2
\end{namedformula}

\begin{importantnote}
    If the image from lens 1 lies beyond lens 2, then the object distance for
    lens 2 is negative. This is the unusual case that deserves extra attention.
\end{importantnote}

\subsection{Worked Example: Two Converging Lenses}

Let \(f_1=\SI{24}{cm}\), \(f_2=\SI{9.0}{cm}\), lens separation
\(L=\SI{10}{cm}\), and \(p_1=\SI{6.0}{cm}\).

For lens 1:

\[
    \frac{1}{6.0}+\frac{1}{i_1}=\frac{1}{24}
\]

\[
    \frac{1}{i_1}
    =
    \frac{1}{24}-\frac{1}{6.0}
    =
    -\frac{1}{8}
\]

\[
    \boxed{i_1=-\SI{8.0}{cm}}
\]

The first image is virtual and lies \(\SI{8.0}{cm}\) to the left of lens 1.
Lens 2 is \(\SI{10}{cm}\) to the right of lens 1, so

\[
    p_2=10+8=\SI{18}{cm}.
\]

For lens 2:

\[
    \frac{1}{18}+\frac{1}{i_2}=\frac{1}{9.0}
\]

\[
    \boxed{i_2=\SI{18}{cm}}
\]

The final image is real and lies \(\SI{18}{cm}\) to the right of lens 2.

\[
    m_1=-\frac{-8}{6}=\frac{4}{3},
    \qquad
    m_2=-\frac{18}{18}=-1,
    \qquad
    \boxed{M=-\frac{4}{3}}.
\]

The final image is inverted relative to the original object.

\section{Optical Instruments}

Optical instruments use mirrors and lenses to change the apparent size, position,
or clarity of images.

\subsection{Simple Magnifying Lens}

A simple magnifier places the object near the focal point of a converging lens so
that the eye sees a larger virtual image.

\begin{namedformula}{Simple magnifier angular magnification}
    m_u=\frac{\SI{25}{cm}}{f}
\end{namedformula}

The \(\SI{25}{cm}\) distance is the standard near-point reference distance used
in this treatment.

\subsection{Compound Microscope}

A compound microscope uses:

\begin{itemize}
    \item an objective lens to form a real enlarged intermediate image;
    \item an eyepiece to magnify that intermediate image angularly.
\end{itemize}

\begin{namedformula}{Objective magnification}
    m\approx-\frac{s}{f_{\text{ob}}}
\end{namedformula}

\begin{namedformula}{Compound microscope magnification}
    M
    =
    mm_u
    \approx
    -
    \frac{s(\SI{25}{cm})}
    {f_{\text{ob}}f_{\text{ey}}}
\end{namedformula}

\subsection{Refracting Telescope}

A refracting telescope uses a large objective lens to gather light from a distant
object and form an intermediate image. The eyepiece then lets the eye view that
image comfortably.

\begin{namedformula}{Refracting telescope angular magnification}
    m_u=-\frac{f_{\text{ob}}}{f_{\text{ey}}}
\end{namedformula}

The negative sign means the final image is inverted.

\section{Problem-Solving Map}

\begin{enumerate}
    \item Identify the device: plane mirror, spherical mirror, refracting
    surface, lens, or instrument.
    \item Decide which sign convention applies.
    \item Decide whether the image should be real or virtual before calculating.
    \item Use the correct formula.
    \item Interpret the sign of \(i\): positive means real, negative means
    virtual.
    \item Interpret the sign of \(m\): positive means upright, negative means
    inverted.
\end{enumerate}

\begin{importantnote}
    In image problems, a sign is not just arithmetic. It tells you where the
    image is and whether it is physically projectable.
\end{importantnote}

\section{Formula Sheet}

\begin{multicols}{2}

\subsection*{Plane Mirror}

\[
    i=-p
\]

\[
    m=+1
\]

\subsection*{Spherical Mirror}

\[
    f=\frac{r}{2}
\]

\[
    \frac{1}{p}+\frac{1}{i}=\frac{1}{f}
\]

\[
    m=-\frac{i}{p}
\]

\subsection*{Refracting Surface}

\[
    \frac{n_1}{p}+\frac{n_2}{i}
    =
    \frac{n_2-n_1}{r}
\]

\subsection*{Thin Lens}

\[
    \frac{1}{p}+\frac{1}{i}=\frac{1}{f}
\]

\[
    \frac{1}{f}
    =
    (n-1)
    \left(
        \frac{1}{r_1}
        -
        \frac{1}{r_2}
    \right)
\]

\[
    m=-\frac{i}{p}
\]

\subsection*{Multiple Lenses}

\[
    M=m_1m_2
\]

\subsection*{Simple Magnifier}

\[
    m_u=\frac{\SI{25}{cm}}{f}
\]

\subsection*{Microscope}

\[
    M\approx
    -
    \frac{s(\SI{25}{cm})}
    {f_{\text{ob}}f_{\text{ey}}}
\]

\subsection*{Telescope}

\[
    m_u=-\frac{f_{\text{ob}}}{f_{\text{ey}}}
\]

\end{multicols}

\section{Common Mistakes}

\begin{enumerate}
    \item Treating a virtual image as though it can be projected onto a screen.
    \item Forgetting that plane mirrors have \(i=-p\).
    \item Using the mirror radius sign convention for refracting surfaces.
    \item Forgetting that convex mirrors produce only virtual upright images.
    \item Forgetting that diverging lenses produce only virtual upright images
    for real objects.
    \item Using \(m=-i/p\) but ignoring the meaning of the sign.
    \item In a two-lens system, measuring \(p_2\) from the wrong place.
    \item Forgetting that telescope magnification is angular, not lateral.
\end{enumerate}

\end{document}
